\documentclass[11pt,a4paper]{article}
\usepackage[utf8]{inputenc}
\usepackage[T1]{fontenc}
\usepackage{amsmath,amssymb,amsthm}
\usepackage{listings}
\usepackage{geometry}
\usepackage{hyperref}
\geometry{margin=2.5cm}

\lstset{
  language=C++,
  basicstyle=\ttfamily\small,
  keywordstyle=\bfseries,
  breaklines=true,
  numbers=left,
  numberstyle=\tiny,
  frame=single,
  tabsize=2
}

\title{IOI 2024 -- Tree}
\author{Solution}
\date{}

\begin{document}
\maketitle

\section{Problem Summary}
Given a rooted tree with $N$ nodes, each with weight $W[i]$. Assign integer ``coefficients'' $C[i]$ to each node such that for every node $i$, the sum of coefficients in its subtree is between $L$ and $R$ (given per query). Minimize the total cost $\sum |C[i]| \cdot W[i]$.

\section{Solution Approach}

\subsection{Key Observation}
Let $S[i]$ denote the sum of coefficients in the subtree of $i$. Then $L \le S[i] \le R$ for all $i$, and $C[i] = S[i] - \sum_{j \in \text{children}(i)} S[j]$.

So the cost is:
\[\sum_{i=0}^{N-1} |S[i] - \sum_{j \in \text{children}(i)} S[j]| \cdot W[i]\]

We want to choose $S[i] \in [L, R]$ for all $i$ to minimize this cost.

\subsection{Greedy / DP on Tree}
Process the tree bottom-up. For each leaf $i$: $S[i] \in [L, R]$ and $C[i] = S[i]$, so cost contribution is $|S[i]| \cdot W[i]$. Since $L \ge 1$ (usually), choose $S[i] = L$ to minimize.

For an internal node $i$ with children $j_1, \ldots, j_k$:
\begin{itemize}
  \item The children have already chosen their $S[j]$ values.
  \item Let $T = \sum_j S[j]$. Then $C[i] = S[i] - T$.
  \item To minimize $|C[i]| \cdot W[i]$, choose $S[i]$ as close to $T$ as possible, subject to $S[i] \in [L, R]$.
  \item So $S[i] = \text{clamp}(T, L, R)$ and $|C[i]| = |S[i] - T|$.
\end{itemize}

\subsection{Algorithm}
\begin{enumerate}
  \item For each leaf: $S[i] = L$.
  \item For each internal node (bottom-up): $T = \sum_{j \in \text{children}} S[j]$. $S[i] = \max(L, \min(R, T))$.
  \item Cost contribution: $|S[i] - T| \cdot W[i]$ for internal nodes, $|S[i]| \cdot W[i] = L \cdot W[i]$ for leaves.
\end{enumerate}

\section{C++ Solution}

\begin{lstlisting}
#include <bits/stdc++.h>
using namespace std;
typedef long long ll;

int N;
vector<int> par, weight;
vector<vector<int>> children;
vector<int> order; // topological order (leaves first)

void init(vector<int> P, vector<int> W) {
    N = P.size();
    par = P;
    weight = W;
    children.resize(N);
    for (int i = 1; i < N; i++)
        children[P[i]].push_back(i);

    // Compute topological order (bottom-up)
    order.clear();
    vector<bool> visited(N, false);
    // BFS from leaves
    queue<int> q;
    vector<int> degree(N, 0);
    for (int i = 0; i < N; i++)
        degree[i] = children[i].size();

    for (int i = 0; i < N; i++)
        if (degree[i] == 0)
            q.push(i);

    while (!q.empty()) {
        int u = q.front(); q.pop();
        order.push_back(u);
        if (par[u] >= 0) {
            degree[par[u]]--;
            if (degree[par[u]] == 0)
                q.push(par[u]);
        }
    }
}

ll query(int L, int R) {
    vector<ll> S(N);
    ll total_cost = 0;

    for (int u : order) {
        if (children[u].empty()) {
            // Leaf
            S[u] = L;
            total_cost += (ll)L * weight[u];
        } else {
            ll child_sum = 0;
            for (int c : children[u])
                child_sum += S[c];

            S[u] = max((ll)L, min((ll)R, child_sum));
            total_cost += abs(S[u] - child_sum) * weight[u];
        }
    }

    return total_cost;
}

int main() {
    int n, q;
    scanf("%d %d", &n, &q);
    vector<int> P(n), W(n);
    P[0] = -1;
    for (int i = 1; i < n; i++) scanf("%d", &P[i]);
    for (int i = 0; i < n; i++) scanf("%d", &W[i]);
    init(P, W);
    while (q--) {
        int L, R;
        scanf("%d %d", &L, &R);
        printf("%lld\n", query(L, R));
    }
    return 0;
}
\end{lstlisting}

\section{Complexity Analysis}
\begin{itemize}
  \item \textbf{Preprocessing:} $O(N)$ for topological sort.
  \item \textbf{Per query:} $O(N)$ for bottom-up traversal.
  \item \textbf{Space:} $O(N)$.
\end{itemize}

\textbf{Optimization:} For queries with $L = 1$, the structure simplifies: leaves always have $S = 1$, and the optimal $S[i]$ values depend on the tree structure. We can precompute the answer as a function of $R$ using the observation that increasing $R$ only affects nodes where the children's sum exceeds $R$ (clamping reduces the excess). With careful analysis, the answer is a piecewise linear function of $R$, enabling $O(\log N)$ per query after $O(N \log N)$ preprocessing.

\end{document}
